% Generated draft: requires mathematical and source review.
\documentclass[11pt]{article}
\usepackage[margin=25mm]{geometry}
\usepackage{fontspec}
\setmainfont{Latin Modern Roman}
\usepackage{amsmath,amssymb,mathtools,unicode-math}
\setmathfont{Latin Modern Math}
\usepackage{enumitem,needspace,xurl,hyperref}
\setlength{\parindent}{0pt}
\setlength{\parskip}{5pt}
\setlength{\emergencystretch}{3em}
\newcommand{\N}{\mathbb{N}}
\newcommand{\Z}{\mathbb{Z}}
\newcommand{\Q}{\mathbb{Q}}
\newcommand{\R}{\mathbb{R}}
\newcommand{\C}{\mathbb{C}}
\newcommand{\F}{\mathbb{F}}
\newcommand{\A}{\mathbb{A}}
\newcommand{\PP}{\mathbb{P}}
\newcommand{\E}{\mathbb{E}}
\newcommand{\eps}{\varepsilon}
\newcommand{\dd}{\,\mathrm{d}}
\begin{document}
\section*{UnsolvedMath ID: 30002666}
\textbf{Source problem code:} catalogue entry
\subsection{Definitions and mathematical statement}
Let \(K\) denote the two-dimensional complex described in the supplied record as being formed from triangular faces by identifying edges and/or vertices, and let \(|K|\) denote its underlying topological space, with the quotient topology induced by these identifications. The faces and identifications specifying the intended \(K\) are not supplied.

In the topological interpretation, an embedding of \(|K|\) in \(\R^3\) is a continuous map \(f:|K|\to\R^3\) that is a homeomorphism onto its image, where \(f(|K|)\) has the subspace topology. The question for the intended complex is whether
\[
\exists f:|K|\longrightarrow\R^3
\quad\text{such that }f\text{ is a topological embedding}.
\]
The source does not specify whether a topological, piecewise-linear, or other embedding category is intended; the displayed formulation makes only the topological interpretation explicit. It also does not state finiteness or local finiteness hypotheses, specify the gluing maps, or identify the particular complex. Consequently, this is an embeddability question with an unspecified input, not a fully specified yes-or-no instance or an assertion that every such complex embeds. The supplied background refers to a pictured complex, but no picture or equivalent combinatorial description is included.
\subsection{Short English statement}
Can the intended triangular two-dimensional complex be embedded in \(\R^3\)? In the topological interpretation, this asks for a homeomorphism onto a subspace of \(\R^3\). The supplied record omits the intended complex and leaves the embedding category unspecified, so the particular question requires clarification.
\subsection{Sources}
\begin{itemize}
\item Local UnsolvedMath registry, supplied complete record for permanent ID \(30002666\), source problem code \texttt{catalogue entry}, titled ``Embedding Triangular 2-Complexes in 3-Space.'' The supplied record labels the problem \texttt{open} and the statement \texttt{corrected\_verified}; these are supplied curation labels, not independent verification here. Its statement-verification note says that notation and publication prose were repaired without changing the substance.
\item Supplied report citation: \emph{Discrete Geometry (2015)}, Oberwolfach Reports, DOI 10.4171/owr/2014/40, pp.~2235--2296. \url{https://doi.org/10.4171/owr/2014/40}. The background describes Oberwolfach Reports as documenting mathematical workshops and their open-problem discussions. The record gives proposed year 2014 but supplies no named proposer.
\item Supplied literature source: ``Sedgwick, Oberwolfach report.'' \url{https://ems.press/content/serial-article-files/46530}. This source label does not by itself establish who proposed the problem.
\item Supplied literature source: ``Matou\v{s}ek--Sedgwick--Tancer--Wagner, embeddability decision theorem.'' \url{https://arxiv.org/abs/1402.0815}.
\item Supplied literature assessment dated exactly \texttt{2026-08-21}: status ``open''; the assessment states that embeddability of arbitrary two-complexes in three-space is decidable, but that the extracted record does not specify the pictured complex, preventing assessment of its particular yes-or-no instance. This is supplied dated curation derived from the OWR material, not an independently verified literature assessment. The supplied background expressly calls for independent verification of time-sensitive status and cited claims.
\end{itemize}
\end{document}
